\section{Principle and Action 23}

\paragraph{Chapter 23}

\begin{verse}
Speak little and act like nature\\ A violent wind does not blow all morning\\ A sudden downpour does not last all day\\ 
This impermanence is proper to nature itself = to Heaven-and-Earth \\ 
If even nature does not have the power to persist (in extreme actions)\\ 
Man still less can do so\\ 
Therefore: in life, imitate (by non-acting) the Principle\\ 
To become one with the Principle\\ 
Imitate (by acting) its Virtue (\emph{tê}) to become one with it\\ 
Abandon yourself to transformations to realize in yourself its power\\ 
By aligning with the Principle, partake of it\\ 
By aligning with the Virtue, partake of it\\ 
By aligning with transformations, partake of them\\ 
(But in most) this inner security is lacking, hence incomplete trust.
\end{verse}

\paragraph{Commentary}


The True Man aligns himself with the Principle and reproduces it in its dual nature: non-acting (the Tao in the proper sense) and acting (the Virtue of the Tao). The latter, as is known, operates in perpetual transformation. Thus, in all that is external, the True Man does not stiffen nor resist the instability that even Heaven-and-Earth, especially in its extreme manifestations, reveals to him; he follows the course of the transformations of forms, the immutable not being something to grasp or realize on this plane.

However, such a path requires inner security, something like intrepidity or ``metaphysical trust'' (in the Way): the art of an active letting-go, like one who abandons himself to the play of the waves without being submerged. This, indeed, is for the few (final line).

In the tenth and thirteenth lines, the character \emph{shih}, rendered as ``transformation,'' more commonly means ``defect'' (``transformation'' would properly be \emph{i}). The literal translation would therefore be: ``to unite with the defect (\emph{shih}) in order to partake of it (or enjoy it).'' With ``transformation'' (referring to the instability and changes discussed at the beginning and which are the mode of the Way) instead of ``defect,'' the meaning of the two lines aligns well with the rest. However, if one retains the literal meaning, the phrase could once again point to the principle of being whole within the fragment and of one's own nature: uniting with the defect is equivalent to embracing what exists within the limiting conditions of the formal state; through this alignment, excluding arbitrariness and the mania of the ``I'', one is in the Way.

\paragraph{Chinese text and literal translation}


Chapter 23 (\begin{CJK*}{UTF8}{bsmi}第二十三章\end{CJK*})}

\columnratio{0.37}
\begin{paracol}{2}
\begin{CJK*}{UTF8}{bsmi}
\begin{verse}
希言自然。\\ 
故飄風不終朝，\\ 
驟雨不終日。\\ 
孰為此者？\\ 
天地尚不能久，\\ 
而況於人乎？\\ 
故從事於道者，\\ 
道者同於道，\\ 
德者同於德，\\ 
失者同於失。\\ 
同於道者，道亦樂得之；\\ 
同於德者，德亦樂得之；\\ 
同於失者，失亦樂得之。\\ 
信不足焉，\\ 
有不信焉。
\end{verse}
\end{CJK*}
\switchcolumn
\vspace*{-2ex}
\begin{verse}
Less spoken, words speak for themselves naturally.\\ 
Therefore gusts cannot chill the exuberant day,\\ 
Showers cannot turn day to dusk.\\ 
Why is it so?\\ 
Even the heaven and earth cannot counter its own force,\\ 
Let alone people?\\ 
Therefore those engaged with the Dao take pleasure that,\\ 
The Taught is the Dao,\\ 
The Virtuous is the Virtue,\\ 
The perplexed is perplexity.\\ 
At one with the Dao, the Dao is welcoming;\\ 
At one with Virtue, Virtue be appreciating;\\ 
At one with perplexity, perplexity be gratifying.\\ 
Lack of belief,\\ 
Explains the disbelief in this futility.
\end{verse}
\end{paracol}

\flrightit{Posted on 2025-05-06 by Tannheuser}

